\section{Axis-Based Representation}
\label{sec:axis-model}
Within the read-once fragment, a computation is determined by which input
axes share a coordinate and which axes remain in the output, in what order.
This section makes that observation precise. All statements use
Section~\ref{sec:formal-scope}'s assumptions.

\subsection{Logical Axes}
A logical axis groups input axes that use the same index. For example, in
\code{einsum('ik,kj->ij', A, B)}, the two axes named \code{k} share an index.
We record this relation using input \emph{ports}. Port $(X,d)$ is dimension
$d$ of input $X$. For a component with input set $\mathcal X$, let
\begin{equation}
 P=\{(X,d)\mid X\in\mathcal X,\ 0\le d<\operatorname{rank}(X)\}.
\end{equation}
Two ports are related by $\sim$ when they are connected, directly or
indirectly, to use the same index. A logical axis is one equivalence class:
\begin{equation}
 \mathcal A=P/{\sim}.
\end{equation}
Equal extents do not establish axis identity. We also retain the classes of
axes removed by summation: they still record how input coordinates connect.

\subsection{The Representation}
For an expression $T$ with assigned attributes, define
\begin{equation}
 N(T)=(P,\sim,F),
 \label{eq:axis-representation}
\end{equation}
where $F$ lists the output classes in dimension order. This representation
specifies the value at each output coordinate: use the same index at
connected input ports, multiply the selected input entries, and sum over
all classes absent from $F$. Classes in $F$ remain as output coordinates.

For matrix multiplication $T=AB$, the classes are
\begin{equation}
 i=\{(A,0)\},\qquad k=\{(A,1),(B,0)\},\qquad j=\{(B,1)\}.
\end{equation}
Thus
\begin{equation}
 N(T)=\bigl(\{(A,0),(A,1),(B,0),(B,1)\},\sim_{AB},(i,j)\bigr).
\end{equation}
The class $k$ is absent from the output, giving $T_{ij}=\sum_k A_{ik}B_{kj}$.

\subsection{When the Representation Determines Equivalence}
\begin{lemma}[Axis exactness]
\label{lem:axis-exactness}
For two legal components $L,R$ over the same inputs at fixed source metadata,
\begin{equation}
 \forall X,\quad L(X)=R(X)\quad\iff\quad N(L)=N(R).
 \label{eq:axis-exactness}
\end{equation}
\end{lemma}
Equality of representations means the same connections between input ports
and the same class at each output position, regardless of the names assigned
to classes. Tensor equality includes the ordered output shape.

We use three properties of read-once trees. Every class contains an input
port, no class contains two ports of the same input, and a logical axis
appears at most once in any node's output. Contraction connects axes from
disjoint input subtrees. If two separate axes survive in the same child,
a later operation cannot merge them with each other. A removed axis cannot
reappear. These properties also support the synthesis algorithm below.

\begin{proof}
If $N(L)=N(R)$, each output coordinate selects the same input entries and
sums over the same index combinations. The resulting sums of products agree.

For the converse, suppose $N(L)\ne N(R)$. Different output shapes already
show inequivalence, so consider equal shapes. Add a port for each output
position and connect it to its class. Output-order differences now appear
as differences in connections too. There are two ports connected in one
representation but separate in the other. Swapping $L,R$ if necessary, call
them $u,w$ with $u\sim_L w$ and $u\not\sim_R w$.

In $R$, assign index 0 to the class of $u$, index 1 to the class of $w$,
and index 0 to every other class. These indices exist because extents are
at least two. Set each input to a one-hot tensor $X^*$ with value 1 only at
the selected coordinate, and choose output coordinate $y$ the same way.
Set rank-zero inputs to 1. Every class has an input port, so exactly one
summand of $R(X^*)[y]$ is 1. In $L$, the connected ports would require
indices 0 and 1 at once, making every summand zero. Therefore
\begin{equation}
 R(X^*)[y]=1\ne0=L(X^*)[y].
\end{equation}
This proves the contrapositive of the converse.
\end{proof}

For example, a dot product differs from multiplying two independent sums.
The one-hot inputs $A^*=(1,0)$ and $B^*=(0,1)$ distinguish them:
\begin{equation}
\begin{aligned}
 L(A^*,B^*)&=\sum_k A^*_k B^*_k=1\cdot0+0\cdot1=0,\\
 R(A^*,B^*)&=\left(\sum_i A^*_i\right)
             \left(\sum_j B^*_j\right)=(1+0)(0+1)=1.
\end{aligned}
\end{equation}
The axis extents are the same; the connections are different.
